\subsection{UNION AND INTERSECTION OF TWO SETS}

If A, B are sets, we write

\[
\mathrm{A} \cup \mathrm{B} = \bigcup_ {\mathbf {X} \in \{\mathbf {A}, \mathrm{B} \}} \mathrm{X}, \quad \mathrm{A} \cap \mathrm{B} = \bigcap_ {\mathbf {X} \in \{\mathbf {A}, \mathrm{B} \}} \mathrm{X}.
\]

It is clear that A ∪ B is the set of all objects which belong either to A or to B (or possibly to both), while A ∩ B is the set of all objects which belong to both A and B. In particular, \{x, y\} = \{x\} ∪ \{y\}. Let \{x, y, z\} = \{x, y\} ∪ \{z\}. The set \{x, y, z\} is the set whose only elements are x, y, and z. Similarly we write

\[
\{x, y, z, t \} = \{x, y, z \} \cup \{t \},
\]

and so on.

If now A, B, C, D are sets, we write

\[
\begin{array}{c} \mathrm {A \cup B \cup C = \bigcup_ {\mathbf {X} \in \{\mathrm{A,B,C} \}} X, \quad A \cap B \cap C = \bigcap_ {\mathbf {X} \in \{\mathrm{A,B,C} \}} X;} \\ \mathrm {A \cup B \cup C \cup D = \bigcup_ {\mathbf {X} \in \{\mathrm{A,B,C,D} \}} X, \quad A \cap B \cap C \cap D = \bigcap_ {\mathbf {X} \in \{\mathrm{A,B,C,D} \}} X;} \end{array}
\]

and so on.

Let A, B, C be sets. From Propositions 1 and 2 we deduce the formulae

\[
\begin{array}{c} \mathbf {A} \cup \mathbf {B} = \mathbf {B} \cup \mathbf {A}, \qquad \mathbf {A} \cap \mathbf {B} = \mathbf {B} \cap \mathbf {A}, \\ \mathbf {A} \cup (\mathbf {B} \cup \mathbf {C}) = (\mathbf {A} \cup \mathbf {B}) \cup \mathbf {C} = \mathbf {A} \cup \mathbf {B} \cup \mathbf {C}, \\ \mathbf {A} \cap (\mathbf {B} \cap \mathbf {C}) = (\mathbf {A} \cap \mathbf {B}) \cap \mathbf {C} = \mathbf {A} \cap \mathbf {B} \cap \mathbf {C}. \end{array}
\]

These formulae are also immediate consequences of the theorems enunciated in the criterion C24 (Chapter I, §3, no. 5). Similarly one proves the formulae

\[
\mathbf {A} \cup (\mathbf {B} \cap \mathbf {C}) = (\mathbf {A} \cup \mathbf {B}) \cap (\mathbf {A} \cup \mathbf {C}), \quad \mathbf {A} \cap (\mathbf {B} \cup \mathbf {C}) = (\mathbf {A} \cap \mathbf {B}) \cup (\mathbf {A} \cap \mathbf {C})
\]

(``distributivity'' of union over intersection and of intersection over union; cf.~§5, no. 6).

The relation \(\mathbf{A} \subset \mathbf{B}\) is equivalent to \(\mathbf{A} \cup \mathbf{B} = \mathbf{B}\) and to \(\mathbf{A} \cap \mathbf{B} = \mathbf{A}\) . If \(\mathbf{A}\) and \(\mathbf{B}\) are subsets of a set \(\mathbf{E}\) , we deduce from Proposition 5 (or from criterion C24) the formulae

\[
\mathbb {G} _ {\mathbf {E}} (\mathbf {A} \cup \mathbf {B}) = \left(\mathbb {G} _ {\mathbf {E}} \mathbf {A}\right) \cap \left(\mathbb {G} _ {\mathbf {E}} \mathbf {B}\right), \quad \mathbb {G} _ {\mathbf {E}} (\mathbf {A} \cap \mathbf {B}) = \left(\mathbb {G} _ {\mathbf {E}} \mathbf {A}\right) \cup \left(\mathbb {G} _ {\mathbf {E}} \mathbf {B}\right);
\]

furthermore, we have

\[
\mathrm{A} \cup (\mathbb {C} _ {\mathrm{E}} \mathrm{A}) = \mathrm{E}, \quad \mathrm{A} \cap (\mathbb {C} _ {\mathrm{E}} \mathrm{A}) = \varnothing .
\]

If \(\Gamma\) is a correspondence between \(E\) and \(F\) , and if \(A, B\) are subsets of \(E\) , it follows from Proposition 3 that

\[
\Gamma \langle \mathrm{A} \cup \mathrm{B} \rangle = \Gamma \langle \mathrm{A} \rangle \cup \Gamma \langle \mathrm{B} \rangle , \quad \Gamma \langle \mathrm{A} \cap \mathrm{B} \rangle \subset \Gamma \langle \mathrm{A} \rangle \cap \Gamma \langle \mathrm{B} \rangle ,
\]

and that, if \(f\) is a mapping of \(F\) into \(E\) ,

\[
\overline {{{f}}} ^ {1} \langle \mathrm{A} \cap \mathrm{B} \rangle = \overline {{{f}}} ^ {1} \langle \mathrm{A} \rangle \cap \overline {{{f}}} ^ {1} \langle \mathrm{B} \rangle
\]

from Proposition 4.

We record also the following Proposition on complements :

PROPOSITION 6. Let \(f\) be a mapping of \(A\) into \(B\) . For every subset \(Y\) of \(B\) , we have

\[
\overline {{{f}}} ^ {1} \langle \mathrm{B} - \mathrm{Y} \rangle = \overline {{{f}}} ^ {1} \langle \mathrm{B} \rangle - \overline {{{f}}} ^ {1} \langle \mathrm{Y} \rangle .
\]

For \(x\) belongs to \(\overline{f}^1\langle \mathbf{B} - \mathbf{Y}\rangle\) if and only if \(f(x)\) belongs to \(\mathbf{B}\) but not to \(\mathbf{Y}\) , i.e., if and only if \(x\) belongs to \(\overline{f}^1\langle \mathbf{B}\rangle\) but not to \(\overline{f}^1\langle \mathbf{Y}\rangle\) .

COROLLARY. Let \(f\) be an injection of \(\mathbf{A}\) into \(\mathbf{B}\) . For every subset \(\mathbf{X}\) of \(\mathbf{A}\) , we have \(f\langle \mathbf{A} - \mathbf{X}\rangle = f\langle \mathbf{A}\rangle - f\langle \mathbf{X}\rangle\) .

Writing \(f = i \circ g\) , where \(i\) is the canonical injection of \(f\langle A\rangle\) into B, we reduce the Corollary to Proposition 6 applied to \(\overline{g}^1\) .

¶ The intersection X ∩ A is sometimes called the trace of X on A. If F is a family of sets, the set of traces on A of the sets belonging to F is called the trace of F on A.

